% !TeX program = pdflatex
\documentclass[aspectratio=169,xcolor={dvipsnames,table}]{beamer}

\input{../../_en_preambulo_beamer}
% Comandos y macros estadísticas compartidas para las presentaciones Beamer en inglés (Metropolis)

% Conjuntos numéricos básicos
\newcommand{\R}{\mathbb{R}}
\newcommand{\N}{\mathbb{N}}
\newcommand{\Q}{\mathbb{Q}}
\newcommand{\Z}{\mathbb{Z}}
\newcommand{\set}[1]{\left\{ #1 \right\}}
\newcommand{\abs}[1]{\left|#1\right|}
\newcommand{\norm}[1]{\left\| #1 \right\|}

% Comandos estadísticos y probabilísticos del curso
\newcommand{\Var}{\operatorname{Var}}
\newcommand{\cov}{\operatorname{Cov}}
\newcommand{\corr}{\rho}
\newcommand{\comb}[2]{\begin{pmatrix} #1 \\ #2 \end{pmatrix}}
\newcommand{\s}{\sigma}
\newcommand{\particion}{\mathfrak{P}}
\newcommand{\card}[1]{\#\left( #1 \right)}
\newcommand{\seq}[3]{\set{#1}_{#2}^{#3}}
\newcommand{\E}{\mathbb{E}}
\newcommand{\Prob}{\mathbb{P}}

% Entornos pedagógicos y cajas en inglés académico natural (soportando tanto nombres en inglés como en español)
\newenvironment{discussion}{%
  \begin{alertblock}{Discussion Question}%
}{%
  \end{alertblock}%
}
\newenvironment{discusion}{%
  \begin{alertblock}{Discussion Question}%
}{%
  \end{alertblock}%
}

\newenvironment{reflection}{%
  \begin{alertblock}{Classroom Reflection}%
}{%
  \end{alertblock}%
}
\newenvironment{reflexion}{%
  \begin{alertblock}{Classroom Reflection}%
}{%
  \end{alertblock}%
}

\newenvironment{paradox}{%
  \begin{alertblock}{Exploratory Paradox}%
}{%
  \end{alertblock}%
}
\newenvironment{paradoja}{%
  \begin{alertblock}{Exploratory Paradox}%
}{%
  \end{alertblock}%
}

\newenvironment{motivation}{%
  \begin{exampleblock}{Motivation: Data Science in Practice}%
}{%
  \end{exampleblock}%
}
\newenvironment{motivacion}{%
  \begin{exampleblock}{Motivation: Data Science in Practice}%
}{%
  \end{exampleblock}%
}

\newenvironment{quickchallenge}{%
  \begin{block}{Quick Team Challenge}%
}{%
  \end{block}%
}
\newenvironment{retoclase}{%
  \begin{block}{Quick Team Challenge}%
}{%
  \end{block}%
}


\title{Normal Distribution}
\subtitle{Section 03.09 --- Continuous Approximation of Discrete Variables and CLT}
\author[J. Castillo Colmenares]{Juliho Castillo Colmenares}
\institute[Tec de Monterrey]{Tecnologico de Monterrey}
\date{\vspace{-1.5cm}}

\begin{document}

% SLIDE 01: Title Page
\begin{frame}[plain]
	\titlepage
\end{frame}

% SLIDE 02: Roadmap
\begin{frame}{Roadmap --- Chapter 03: Discrete Random Variables}
	\vspace{-0.15cm}
	\begin{block}{Curricular Structure of Unit 2}
		\footnotesize
		\begin{enumerate}\itemsep=0.03cm
			\item \textcolor{gray}{Section 03.01 to 03.03: Foundations, PMF, CDF, Expectation, and Variance}
			\item \textcolor{gray}{Section 03.04 to 03.06: Binomial, Geometric, Negative Binomial, Hypergeometric}
			\item \textcolor{gray}{Section 03.07: Poisson Distribution}
			\item \textcolor{gray}{Section 03.08: Multinomial Distribution}
			\item \textbf{\textcolor{TecRojo}{Section 03.09: Normal Distribution}} $\leftarrow$ \emph{Current focus}
		\end{enumerate}
	\end{block}
	\vspace{-0.2cm}
	\begin{alertblock}{Learning Objective}
		\scriptsize
		Introduce the Normal distribution as the fundamental continuous model, master Z-score standardization and the 68-95-99.7 rule, exploit the De Moivre-Laplace theorem to approximate the Binomial by Normal with continuity correction, and apply the Central Limit Theorem to sums of discrete variables.
	\end{alertblock}
\end{frame}

% SLIDE 03: Formal Definition and Z-Score
\begin{frame}{Formal Definition and Z-Score Standardization}
	\vspace{-0.15cm}
	\begin{columns}[T]
		\begin{column}{0.48\textwidth}
			\begin{block}{Probability Density Function (PDF)}
				\footnotesize
				A continuous random variable $X \sim N(\mu, \sigma^{2})$ has PDF:
				$$f(x) = \dfrac{1}{\sigma\sqrt{2\pi}} \exp\left(-\dfrac{(x-\mu)^{2}}{2\sigma^{2}}\right)$$
				Bell-shaped curve centered at $\mu$ with spread controlled by $\sigma$.
			\end{block}
		\end{column}
		\begin{column}{0.48\textwidth}
			\begin{alertblock}{Z-Score Standardization}
				\footnotesize
				$Z = (X - \mu)/\sigma \sim N(0, 1)$. Probabilities are computed via the standard Normal CDF $\Phi(z) = P(Z \le z)$ with the symmetry $\Phi(-z) = 1 - \Phi(z)$.
			\end{alertblock}
		\end{column}
	\end{columns}
\end{frame}

% SLIDE 04: De Moivre-Laplace and Yates Correction
\begin{frame}{De Moivre-Laplace and Yates Continuity Correction}
	\vspace{-0.15cm}
	\begin{columns}[T]
		\begin{column}{0.48\textwidth}
			\begin{block}{Binomial-to-Normal Limit}
				\footnotesize
				For $X \sim \text{Bin}(n, p)$ with $np \ge 5$ and $n(1-p) \ge 5$:
				$$\dfrac{X - np}{\sqrt{np(1-p)}} \xrightarrow{d} N(0, 1) \quad \text{as } n \to \infty$$
			\end{block}
		\end{column}
		\begin{column}{0.48\textwidth}
			\begin{alertblock}{Yates Continuity Correction}
				\footnotesize
				To approximate $P(X \ge k)$, use $P(Y \ge k - 0.5)$ where $Y \sim N(np, np(1-p))$. The shift by $0.5$ reduces systematic discretization error and typically halves the approximation error.
			\end{alertblock}
		\end{column}
	\end{columns}
\end{frame}

% SLIDE 05: Poisson-Normal and CLT
\begin{frame}{Poisson-to-Normal Approximation and Central Limit Theorem}
	\vspace{-0.15cm}
	\begin{columns}[T]
		\begin{column}{0.48\textwidth}
			\begin{block}{Poisson-to-Normal Limit}
				\footnotesize
				For $X \sim \text{Pois}(\lambda)$ with $\lambda \ge 30$:
				$$\dfrac{X - \lambda}{\sqrt{\lambda}} \xrightarrow{d} N(0, 1) \quad \text{as } \lambda \to \infty$$
				Useful in queueing theory and risk management.
			\end{block}
		\end{column}
		\begin{column}{0.48\textwidth}
			\begin{alertblock}{Central Limit Theorem (CLT)}
				\footnotesize
				For i.i.d. $X_{1}, \dots, X_{n}$ with mean $\mu$ and variance $\sigma^{2}$:
				$Z_{n} = (\sum X_{i} - n\mu) / (\sigma\sqrt{n}) \xrightarrow{d} N(0, 1)$.
				For $n \ge 30$, the approximation is typically excellent.
			\end{alertblock}
		\end{column}
	\end{columns}
\end{frame}

% SLIDE 06: Applications
\begin{frame}{Applications in Engineering, Science, and Data Science}
	\vspace{-0.1cm}
	\begin{itemize}\itemsep=0.1cm
		\item \textbf{Industrial Quality Control:} Monitoring manufacturing processes with $\bar{x}$ and $R$ control charts based on $N(\mu, \sigma^{2})$. \pause
		\item \textbf{Statistical Inference:} Foundation of $Z$-tests, confidence intervals, ANOVA, and linear regression. \pause
		\item \textbf{Physics and Instrumentation:} Measurement errors, Gaussian noise in communication systems, and uncertainty propagation. \pause
		\item \textbf{Finance and Risk Analysis:} Black-Scholes model for option pricing and Value-at-Risk (VaR) measures.
	\end{itemize}
\end{frame}

% SLIDE 07: Python Lab Block 1
\begin{frame}[fragile]{Python Lab: Normal PDF and Z-Score}
	\vspace{-0.05cm}
	\scriptsize
	Computational verification in SciPy (\texttt{03.09\_normal\_approximation.py}) of the Normal PDF, normalization, and Z-score standardization:
	\vspace{0.02cm}
	{\lstset{basicstyle=\fontsize{5pt}{6pt}\selectfont\ttfamily}
	\lstinputlisting[firstline=15, lastline=37]{../../code/03_variables_aleatorias_discretas/03.09_normal_approximation.py}}
\end{frame}

% SLIDE 08: Python Lab Block 2
\begin{frame}[fragile]{Python Lab: Binomial-to-Normal Approximation}
	\vspace{-0.05cm}
	\scriptsize
	De Moivre-Laplace: comparison of Binomial PMF with Normal PDF using Yates continuity correction:
	\vspace{0.02cm}
	{\lstset{basicstyle=\fontsize{4.5pt}{5.5pt}\selectfont\ttfamily}
	\lstinputlisting[firstline=40, lastline=64]{../../code/03_variables_aleatorias_discretas/03.09_normal_approximation.py}}
\end{frame}

% SLIDE 09: Python Lab Block 3
\begin{frame}[fragile]{Python Lab: Poisson-Normal and CLT}
	\vspace{-0.05cm}
	\scriptsize
	Poisson-to-Normal approximation and Monte Carlo verification of the CLT for sum of 50 exponentials:
	\vspace{0.02cm}
	{\lstset{basicstyle=\fontsize{4.5pt}{5.5pt}\selectfont\ttfamily}
	\lstinputlisting[firstline=73, lastline=100]{../../code/03_variables_aleatorias_discretas/03.09_normal_approximation.py}}
\end{frame}

% SLIDE 10: Python Lab Terminal Output
\begin{frame}[fragile]{Python Lab: Terminal Output}
	\vspace{-0.1cm}
	\begin{block}{}
		\fontsize{4pt}{5pt}\selectfont
		\begin{verbatim}
=== Block 1: Normal PDF & Z-Score Standardization ===
X ~ N(2.50, 0.05) | PDF integral: 1.000000
P(2.42 <= X <= 2.58) exact = 0.890401 | 2*Phi(1.60)-1 = 0.890401
Empirical rule: P(|Z|<=1)=0.6827 | P(|Z|<=2)=0.9545 | P(|Z|<=3)=0.9973

=== Block 2: Binomial-to-Normal Approximation ===
Bin(200, 0.25): mu=50.0, sigma=6.1237
  P(X>=60) exact Binomial:  0.062472
  P(X>=60) with Yates:     0.060410 (Z=1.5513)
  P(X>=60) without Yates:  0.051235 (Z=1.6330)
  Error with Yates:    0.002063
  Error without Yates: 0.011237
Convergence: n=50: exact=0.556138 | normal=0.556231

=== Block 3: Poisson-to-Normal & CLT ===
Pois(100) -> N(100, 100): P(85<=X<=115) exact=0.879276
  Normal approx: 0.878858 | Error: 0.000417
Sum of 100 Pois(2) -> S_100 ~ Pois(200):
  P(S_100<=220) exact=0.924699 | normal=0.926411
Sum of 50 Exp(1) ~ Gamma(50, 1):
  P(45<=S_50<=55) exact=0.520993 | normal=0.520500
Monte Carlo (N=250,000): mean=49.9981 | std=7.0764
  P(45<=S<=55) empirical=0.519140
		\end{verbatim}
	\end{block}
\end{frame}

% SLIDE 11: Exercise Level 1 (Statement)
\begin{frame}{In-Class Exercise --- Level 1: Fundamental (1/2)}
	\vspace{-0.1cm}
	\begin{block}{Problem 3.9.1 --- Normal IQ Distribution (Statement)}
		A continuous random variable $X$ follows a Normal distribution with mean $\mu = 100$ and standard deviation $\sigma = 15$. Specify the PDF $f(x)$ and calculate $P(X \le 115)$ and $P(X \ge 85)$ using Z-score standardization.
	\end{block}
	\vspace{0.15cm}
	\pause
	\begin{alertblock}{Model Setup}
		\begin{itemize}\itemsep=0.04cm
			\item $X \sim N(100, 15^{2})$ with PDF $f(x) = \dfrac{1}{15\sqrt{2\pi}} \exp\left(-\dfrac{(x-100)^{2}}{450}\right)$.
			\item Z-score: $Z = (X - 100)/15 \sim N(0, 1)$.
			\item Use \texttt{scipy.stats.norm.cdf} to evaluate probabilities.
		\end{itemize}
	\end{alertblock}
\end{frame}

% SLIDE 12: Exercise Level 1 (Solution)
\begin{frame}{In-Class Exercise --- Level 1: Fundamental (2/2)}
	\vspace{-0.05cm}
	\scriptsize
	Standardize $Z = (X - 100)/15$ and use the symmetry of the standard Normal: \pause
	\begin{block}{Probability Calculation}
		\begin{itemize}\itemsep=0.05cm
			\item \textbf{Upper tail:} $P(X \le 115) = P\left(Z \le \dfrac{115-100}{15}\right) = P(Z \le 1) = \Phi(1) \approx 0.8413$ \pause
			\item \textbf{Lower tail:} $P(X \ge 85) = P\left(Z \ge \dfrac{85-100}{15}\right) = P(Z \ge -1) = 1 - \Phi(-1) = \Phi(1) \approx 0.8413$
		\end{itemize}
	\end{block}
	\vspace{0.05cm}
	\pause
	\begin{alertblock}{Standard Normal Symmetry}
		\scriptsize
		By bell-curve symmetry, $P(X \le 115) = P(X \ge 85) \approx 0.8413$ ($84.13\%$). The interval $[85, 115]$ contains exactly $\pm 1$ standard deviation from the mean, capturing $\approx 68\%$ of central mass; each external tail is $\approx 16\%$.
	\end{alertblock}
\end{frame}

% SLIDE 13: Exercise Level 2 (Statement)
\begin{frame}{In-Class Exercise --- Level 2: Operational (1/2)}
	\vspace{-0.1cm}
	\begin{block}{Problem 3.9.4 --- Binomial-Normal Approximation in Standardized Exam (Statement)}
		In a standardized exam with $n = 200$ multiple-choice questions with 4 options each, a student answers at random with $p = 0.25$. Use the Normal approximation with Yates correction to calculate $P(X \ge 60)$.
	\end{block}
	\vspace{0.15cm}
	\pause
	\begin{alertblock}{Comparison with Exact Value}
		\scriptsize
		Calculate the Normal approximation and compare with the exact Binomial value. Operational rules: $np = 50 \ge 5$ and $n(1-p) = 150 \ge 5$.
	\end{alertblock}
\end{frame}

% SLIDE 14: Exercise Level 2 (Solution)
\begin{frame}{In-Class Exercise --- Level 2: Operational (2/2)}
	\vspace{-0.05cm}
	\scriptsize
	With $\mu = 50$ and $\sigma = \sqrt{37.5} \approx 6.124$: \pause
	\begin{block}{Yates Approximation}
		\scriptsize
		$P(X \ge 60) \approx P(Y \ge 59.5)$ where $Y \sim N(50, 6.124^{2})$:
		$$P(X \ge 60) \approx P\left(Z \ge \dfrac{59.5 - 50}{6.124}\right) = P(Z \ge 1.551) = 1 - \Phi(1.551) \approx 0.0604$$
	\end{block}
	\vspace{0.05cm}
	\pause
	\begin{alertblock}{Exact Comparison}
		\scriptsize
		Exact Binomial: $P(X \ge 60) = 1 - \text{stats.binom.cdf}(59, 200, 0.25) \approx 0.0625$. Absolute error: $\abs{0.0604 - 0.0625} \approx 0.0021$, confirming the precision of the Yates correction.
	\end{alertblock}
\end{frame}

% SLIDE 15: Exercise Level 3 (Statement)
\begin{frame}{In-Class Exercise --- Level 3: Analytical (1/2)}
	\vspace{-0.1cm}
	\begin{block}{Problem 3.9.7 --- De Moivre-Laplace Theorem Derivation (Statement)}
		Let $X \sim \text{Bin}(n, p)$. Prove that $Z_{n} = (X - np)/\sqrt{np(1-p)}$ converges in distribution to the standard Normal as $n \to \infty$.
	\end{block}
	\vspace{0.1cm}
	\pause
	\begin{alertblock}{Deduction Strategy}
		\scriptsize
		Use the Binomial MGF $M_{X}(t) = (1 - p + p e^{t})^{n}$ and derive the MGF of $Z_{n}$. Apply the Taylor expansion $\log(1 + x) = x - x^{2}/2 + O(x^{3})$ to obtain $M_{Z}(t) = e^{t^{2}/2}$.
	\end{alertblock}
\end{frame}

% SLIDE 16: Exercise Level 3 (Solution)
\begin{frame}{In-Class Exercise --- Level 3: Analytical (2/2)}
	\vspace{-0.05cm}
	\scriptsize
	Derive the MGF of $Z_{n}$ and take the asymptotic limit: \pause
	\begin{block}{MGF of $Z_{n}$ and Asymptotic Limit}
		\scriptsize
		$M_{Z_{n}}(t) = e^{-npt/\sigma_{n}} (1 - p + p e^{t/\sigma_{n}})^{n}$, with $\sigma_{n} = \sqrt{np(1-p)}$.
	\end{block}
	\vspace{0.05cm}
	\pause
	\begin{alertblock}{Taylor Expansion of the Logarithm}
		\scriptsize
		With $e^{t/\sigma_{n}} - 1 = t/\sigma_{n} + t^{2}/(2\sigma_{n}^{2}) + O(n^{-3/2})$:
		\begin{align*}
			\log(1 - p + p e^{t/\sigma_{n}}) &= \log(1 + p(e^{t/\sigma_{n}} - 1)) = p \cdot \dfrac{t}{\sigma_{n}} + \dfrac{p(1-p)t^{2}}{2\sigma_{n}^{2}} + O(n^{-3/2})
		\end{align*}
		Multiplying by $n$ and centering: $n \log(\cdots) - npt/\sigma_{n} \to t^{2}/2$. Therefore:
		$$\lim_{n \to \infty} M_{Z_{n}}(t) = e^{t^{2}/2} \quad \blacksquare$$
	\end{alertblock}
\end{frame}

% SLIDE 17: Exercise Level 4 (Statement)
\begin{frame}{In-Class Exercise --- Level 4: Challenging (1/2)}
	\vspace{-0.1cm}
	\begin{block}{Problem 3.9.9 --- CLT for Sum of Poisson (Statement)}
		Let $X_{1}, \dots, X_{n}$ be i.i.d. Poisson($\lambda$). By CLT, $S_{n}$ converges to $N(n\lambda, n\lambda)$. For $n = 100$ and $\lambda = 2$:
		\begin{itemize}\itemsep=0.04cm
			\item Calculate $P(S_{100} \le 220)$ using the exact Poisson.
			\item Calculate the Normal approximation with continuity correction.
		\end{itemize}
	\end{block}
	\vspace{0.15cm}
	\pause
	\begin{alertblock}{Model Setup}
		\scriptsize
		$S_{100} \sim \text{Pois}(200)$ with $\mu = 200$, $\sigma = \sqrt{200} \approx 14.142$. Use the operational rule $n\lambda = 200 \ge 30$.
	\end{alertblock}
\end{frame}

% SLIDE 18: Exercise Level 4 (Solution)
\begin{frame}{In-Class Exercise --- Level 4: Challenging (2/2)}
	\vspace{-0.05cm}
	\scriptsize
	Compare the exact Poisson with the Normal approximation with Yates correction: \pause
	\begin{block}{Exact Poisson Value}
		\scriptsize
		$P(S_{100} \le 220) = \text{stats.poisson.cdf}(220, mu=200) \approx 0.9247$
	\end{block}
	\vspace{0.05cm}
	\pause
	\begin{alertblock}{Normal Approximation with Yates}
		\scriptsize
		With $Y \sim N(200, 200)$ and $\sigma = \sqrt{200} \approx 14.142$:
		\begin{align*}
			P(S_{100} \le 220) &\approx P\left(Y \le 220.5\right) = \Phi\left(\dfrac{220.5 - 200}{14.142}\right) = \Phi(1.450) \approx 0.9264
		\end{align*}
		The absolute error is $\abs{0.9247 - 0.9264} \approx 0.0017$, confirming the precision of the CLT for the sum of 100 Poisson(2) random variables. \qedhere
	\end{alertblock}
\end{frame}

% SLIDE 19: Closing and Next Steps
\begin{frame}{Closing Section and Modular Perspective}
	\vspace{-0.15cm}
	\begin{block}{Synthesis on the Bell Curve and CLT}
		\scriptsize
		The Normal distribution emerges as the fundamental continuous model that bridges discrete and continuous worlds. From the De Moivre-Laplace theorem to the Central Limit Theorem, the bell curve appears as the universal asymptotic destination of count distributions when the scale parameters grow sufficiently large.
	\end{block}
	\vspace{0.05cm}
	\pause
	\begin{alertblock}{Key Lessons and Next Step}
		\begin{itemize}\itemsep=0.05cm
			\item \textbf{Universality:} The Normal is the limit of Binomial, Poisson, and sum of i.i.d. random variables with finite variance.
			\item \textbf{Z-Score:} Any computation reduces to $\Phi(\cdot)$ via $Z = (X - \mu)/\sigma$.
			\item \textbf{Yates Correction:} Essential to reduce systematic error when approximating discrete distributions by Normal.
			\item \textbf{Next Session:} \textcolor{TecRojo}{Discrete Distributions in Data Science --- MLE Fitting and Overdispersion}.
		\end{itemize}
	\end{alertblock}
\end{frame}

\end{document}
